\section{Related Work}

\subsection{Learner Assessment}

\textbf{Multimodal and process-aware assessment provide complementary evidence about learner performance.} KidsArtBench combines children's artworks with multidimensional teacher ratings and comments, supporting artistic assessment and formative feedback \citep{ye2026multimodal}. SketchMind analyzes conceptual relationships in scientific drawings through sketch reasoning graphs \citep{latif2025sketchmind}, while StatusKT extracts proficiency indicators from authentic mathematical solution processes \citep{park2026statuskt}. These approaches connect observable work to explicit assessment targets, including artistic performance, conceptual understanding, and demonstrated proficiency. Assessment also informs ongoing creative practice. ArtKrit structures reference-based drawing around composition, value, and color, providing computational guidance and feedback during creation rather than only evaluating completed work \citep{ma2025artkrit}.

\textbf{Knowledge tracing complements these approaches by modeling performance across successive interactions.} Deep Knowledge Tracing uses recurrent networks to represent learning histories \citep{piech2015dkt}, while AKT and simpleKT incorporate question characteristics and relationships between current and previous interactions \citep{ghosh2020akt,liu2023simplekt}. Their principal supervision comes from subsequent response correctness, allowing repeated observations to inform estimates of learner state rather than treating each performance independently.

\textbf{The use of assessment for instruction requires relating observed performance to learning goals and subsequent action} \citep{hattie2007feedback,shute2008feedback}. LongTutor examines this relationship directly by separately evaluating historical evidence acquisition, knowledge-state diagnosis, and adaptive teaching. Its results show that strong evidence acquisition does not consistently translate into effective diagnosis or teaching, even when substantial learning histories are available \citep{li2026longtutor}. Similarly, accurate diagnosis in propositional-logic tutoring does not consistently produce actionable feedback \citep{yasir2026confirming}.

\textbf{These findings distinguish access to learner information from its instructional use.} A remaining challenge is determining which assessed characteristics should guide assistance, making the selection of instructionally relevant personal information an important objective for learner representation learning.

\subsection{Personalized Representations}

\textbf{Personalization methods organize historical information through retrieval, structured profiles, and learned embeddings.} LaMP retrieves personal records to condition current predictions \citep{salemi2024lamp}, while HYDRA combines history reranking with shared and user-specific model components \citep{zhuang2024hydra}. User-LLM encodes behavioral histories as continuous representations \citep{ning2025userllm}, and RPM constructs personal factors and retrieves reasoning examples relevant to the current input \citep{kim2026rpm}. These approaches provide different mechanisms for using individual histories without reproducing every record in the model's context.

\textbf{A complementary approach learns what to retain through downstream supervision.} PLUS jointly optimizes a summarizer and reward model so that personal summaries improve preference prediction \citep{nam2026summaries}. Related work jointly learns user models and personalized language models from individual feedback, extending personalization to the generation objective itself \citep{li2024personalized}. In educational dialogue, \citet{duan2026history} optimize student profiles against authentic subsequent responses and train simulators to distinguish matched from mismatched personal information. Their representations capture knowledge, misconceptions, and interaction patterns, demonstrating that learned profiles can preserve more than linguistic style.

\textbf{Personal information can also determine how outputs are evaluated.} P-Check learns to generate dynamic evaluation checklists and weights criteria according to their ability to distinguish individual preferences \citep{seo2026pcheck}. Its evaluation covers both reward prediction and downstream personalized generation. This advances beyond retaining relevant history to learning the standards by which a response should be judged, providing a direct methodological connection between personal information and output selection.

These methods establish that both personal representations and evaluation criteria can be learned through downstream supervision. \textbf{In educational applications, however, successful preference or response prediction does not independently establish why a learner characteristic justifies a particular form of assistance.} The resulting question is how pedagogical judgments can guide the selection of personal information while preserving its support in actual learning records.

\subsection{Student Simulation and Instructional Support}

\textbf{Instructional systems already connect learner information to teaching decisions.} \citet{rafferty2016pomdp} formulate teaching as planning under uncertainty about learner knowledge, while TRAVER combines knowledge tracing with turn-level verification to guide coding-task completion \citep{wang2025traver}. LearnLens integrates error analysis, topic-linked memory, curriculum-grounded generation, and educator oversight to provide personalized science feedback \citep{zhao2025learnlens}. These approaches establish different mechanisms for translating learner information into assistance, with instructional objectives and domain knowledge shaping how that information is used.

\textbf{Evidence-driven frameworks make the rationale for adaptation more explicit.} The EDF framework, presented at AIED 2026, organizes scaffolding around evidence-based inference, pedagogical decisions, and feedback, with its Copa implementation examined in authentic classroom activity \citep{cohn2026edf}. RegulaTutor similarly connects historical evidence to learner interpretations and instructional policies that retain links to supporting records. Its PairRoute procedure uses execution traces to refine the system workflow and construct stage-level model supervision, with the reported refinement focused on historical evidence acquisition \citep{li2026regulatutor}. Evidence-linked teaching is therefore supported by both system design and targeted model improvement, rather than remaining solely a conceptual connection.

\textbf{Student simulation offers another route by modeling the consequences of assistance.} StudentSim combines pooled training with individual specialization to capture learner behavior and responsiveness to guidance \citep{yang2026simulator}. CPT-Agent separates feedback comprehension from writing proficiency and models selective uptake, constrained revision, and skill accumulation \citep{chen2026cptagent}. PROF uses simulated student revisions to optimize feedback generation directly, rather than relying only on predefined feedback attributes \citep{nair2024prof}. Together, these approaches provide behavioral objectives for evaluating and improving support, although conclusions based on simulated responses remain distinct from evidence obtained through actual learner activity.

\textbf{Expert judgments provide complementary supervision, but their adaptation target must be specified.} PERSA uses demonstrations and pairwise preferences to align programming feedback with an instructor's style while preserving diagnostic correctness \citep{kumar2026persa}. Such supervision addresses consistency with a feedback provider, rather than necessarily establishing suitability for a particular learner. Context-sensitive evaluation offers a more direct test of learner adaptation: \citet{borchers2025adaptivity} vary available learner information and examine both output adaptivity and pedagogical quality. This separates whether a system produces acceptable guidance from whether its guidance responds appropriately to individual learning circumstances.

Existing systems therefore already connect learning evidence, personal models, and instructional action. \textbf{A more specific open question is how limited process histories and pedagogical judgments can jointly train the selection of reliable individual characteristics, beyond what preference learning, behavioral simulation, and evidence-driven tutoring already achieve.} Establishing this contribution requires testing both the evidential support for those characteristics and their effect on instructional appropriateness, separately from improvements in general feedback quality.
